\documentclass[a4paper, 10pt, ]{article}

\input{../../KUT000/preamble_v1.1.tex}

% -----------------------------------------------------------------------------

\def\oznacenieCelku{Kolekcia učebných textov}

% -----------------------------------------------------------------------------


\def\KUTporadoveCislo{021}

% \def\oznacenieVerzie{v1.0}
\def\oznacenieVerzie{\phantom{v1.0}}

\def\mesiacRok{marec 2026}

\def\authorslabel{MT}






% -----------------------------------------------------------------------------

\begin{document} 

% -----------------------------------------------------------------------------
% Uvodny nadpis

\noindent
\parbox[t][18mm][c]{0.3\textwidth}{%
\raisebox{-0.9\height}{%
\phantom{.}
\includegraphics[height=18mm]{../../KUT000/URKFEIlogo_v0.1.pdf}%
}%
}%
\parbox[t][18mm][c]{0.7\textwidth}{%
\raggedleft

\sffamily
\fontsize{16pt}{18pt}
\fontseries{sbc}
\selectfont

\noindent
\textcolor[rgb]{0.75, 0.75, 0.75}{\textls[25]{\oznacenieCelku}}
}%

\noindent
\parbox[t][16mm][b]{0.5\textwidth}{%
\raggedright

\color{Gray}
\sffamily

\fontsize{12pt}{12pt}
\selectfont
\mesiacRok

\fontsize{6pt}{10pt}
\selectfont
github.com/OkoliePracovnehoBodu/KUT

\fontsize{8pt}{10pt}
\selectfont
\authorslabel




}%
\parbox[t][16mm][b]{0.5\textwidth}{%
\raggedleft

\sffamily

\fontsize{6pt}{6pt}
\selectfont

\textcolor[rgb]{0.68, 0.68, 0.68}{\oznacenieVerzie}


\fontsize{14pt}{14pt}
\selectfont

\bfseries


\includegraphics[height=12pt]{../../KUT000/KUT_logo_v0.1.pdf}%
{%
\textls[-50]{\KUTporadoveCislo}
}%
}%

% -----------------------------------------------------------------------------



\vspace{6mm}

% ---------------------------------------------
\sffamily
\bfseries
\fontsize{18pt}{21pt}
\selectfont

\begin{flushleft}
    AeroShield: hmotnosť a poloha ťažiska ramena kyvadla
\end{flushleft}

\bigskip

% -----------------------------------------------------------------------------
\normalsize
\normalfont
% -----------------------------------------------------------------------------

\lstset{style=mystyle}










\noindent
\lettrine[lines=1, nindent=1pt, loversize=0.0]{C}{ieľom} 
textu je uvedenie merania hmotnosti ramena kyvadla zariadenia AeroShield a polohy ťažiska tohto ramena. Pre širší kontext sa odporúča dokument \href{run:../../KUT020/TeX/KUT020.pdf}{\textsf{KUT020}}, kde je laboratórne zariadenie AeroShield podrobnejšie opísané.


\bigskip

\noindent
Meranie bolo realizované na zariadení AeroShield označenom číslom 1 (doska plošných spojov označená číslom 1). 


\paragraph{Hmotnosť}

Postup merania: samotné rameno kyvadla odpojené od hriadeľa sa umiestni na váhu pričom sa dbá na mimimalizáciu vplyvu prívodných vodičov motora (rameno na nich nevisí). Zistená hmotnosť ramena kyvadla je $11,57$ [g].


\begin{center}

    % \vspace{-6pt}

    \vbox{%


        \makebox[\textwidth][c]{%
            \includegraphics[width=15cm]{hmotnost_01}%
        }

        \figcaption{
            Meranie hmotnosti ramena kyvadla zariadenia AeroShield. Zistená hmotnosť je $11,57$ [g].
        }
        \label{hmotnost_01}
    }%vbox

    % \vspace{-6pt}

\end{center}


\paragraph{Poloha ťažiska}

Postup odhadu polohy ťažiska: samotné rameno kyvadla odpojené od hriadeľa sa umiestni na horizontálnu tyč kruhového prierezu. Hľadá sa také miesto dotyku ramena s tyčou, pri ktorom je rameno v rovnováhe (neotáča sa ani na jednu stranu) a dbá sa na minimálny vplyv prívodných vodičov motora. Výsledok je na nasledujúcom obrázku.


\begin{center}

    % \vspace{-6pt}

    \vbox{%


        \makebox[\textwidth][c]{%
            \includegraphics[width=15cm]{tazisko_01}%
        }

        \figcaption{
            Rameno kyvadla v rovnovážnej polohe na horizontálnej tyči kruhového prierezu.
        }
        \label{tazisko_01}
    }%vbox

    % \vspace{-6pt}

\end{center}


\noindent
K ramenu sa priloží pravítko a odmeria sa vzdialenosť od osi otáčania (hriadeľa) po miesto dotyku ramena s tyčou. Zistená vzdialenosť je $7,3$ [cm].


\begin{center}

    % \vspace{-6pt}

    \vbox{%


        \makebox[\textwidth][c]{%
            \includegraphics[width=15cm]{tazisko_02}%
        }

        \figcaption{
            Pohľad zhora na rameno kyvadla v rovnovážnej polohe pre odčítanie vzdialenosti medzi osou otáčania (hriadeľom) a miestom dotyku ramena s tyčou. Zistená vzdialenosť je $7,3$ [cm].
        }
        \label{tazisko_02}
    }%vbox

    % \vspace{-6pt}

\end{center}








% ## Meranie (a fotky) hmotnosti a polohy ťažiska ramena kyvadla zariadenia AeroShield.


% Meranie pre AeroShield číslo 1 (PCB označená číslom 1).



% ### Hmotnosť

% Postup merania: samotné rameno kyvadla odpojené od hriadeľa sa umiestni na váhu pričom sa dbá na mimimalizáciu vplyvu prívodných vodičov motora (rameno na nich nevisí). Zistená hmotnosť ramena kyvadla je 11.57 [g].




% ### Poloha ťažiska

% Postup odhadu polohy ťažiska: samotné rameno kyvadla odpojené od hriadeľa sa umiestni na horizontálnu tyč kruhového prierezu. Hľadá sa také miesto dotyku ramena s tyčou, pri ktorom je rameno v rovnováhe (neotáča sa ani na jednu stranu) a dbá sa na minimálny vplyv prívodných vodičov motora. Výsledok je na obrázku IMG_20260304_102538.jpg. K ramenu sa priloží pravítko a odmeria sa vzdialenosť od osi otáčania (hriadeľa) po miesto dotyku ramena s tyčou. Zistená vzdialenosť je 7.3 [cm].


% ![foto](./IMG_20260304_101256.jpg)
% ![foto](./IMG_20260304_102538.jpg)
% ![foto](./IMG_20260304_102549.jpg)
% ![foto](./IMG_20260304_102553.jpg)












% \section{Pohybová rovnica}


% -----------------------------------------------------------------------------

\end{document}

% -----------------------------------------------------------------------------